\chapter{Korzystanie ze źródeł}
\label{chap:zrodla}

Każda informacja, wzór, rysunek, tabela lub dane zaczerpnięte z~cudzego
opracowania muszą zostać opatrzone odwołaniem do źródła. Opisy
bibliograficzne przechowuje się w~pliku \texttt{literatura.bib}, a~wykaz
literatury na końcu pracy tworzy się automatycznie -- trafiają do niego
wyłącznie pozycje przywołane w~tekście, w~kolejności pierwszego cytowania.

\section{Sposób cytowania}
\label{sec:cytowanie}

Odwołanie ma postać numeru w~nawiasie kwadratowym (polecenie
\verb|\cite{klucz}|). Umieszcza się je \textbf{przed} kropką kończącą
zdanie:
\begin{itemize}
  \item \emph{Pompy ciepła mogą w~znacznym stopniu ograniczyć zużycie gazu
        ziemnego do ogrzewania budynków \cite{iea2022}.} -- poprawnie,
  \item \emph{Pompy ciepła mogą w~znacznym stopniu ograniczyć zużycie gazu
        ziemnego do ogrzewania budynków. \cite{iea2022}} -- niepoprawnie.
\end{itemize}

Inne często potrzebne warianty:
\begin{itemize}
  \item kilka źródeł naraz -- \verb|\cite{lund2014,werner2017,iea2022}|
        daje \cite{lund2014,werner2017,iea2022}; kolejne numery są
        automatycznie łączone w~zakres,
  \item wskazanie strony lub rozdziału -- \verb|\cite[s.~25]{duffie2013}|
        daje \cite[s.~25]{duffie2013},
  \item nazwisko autora w~zdaniu -- \verb|\textcite{werner2017}|:
        \emph{\textcite{werner2017} przedstawił przegląd \dots};
        przy więcej niż dwóch autorach pojawia się skrót „i~in.”:
        \emph{\textcite{lund2014} zdefiniowali \dots}
\end{itemize}

Jeżeli promotor wymaga stylu autor--rok, np.\ \emph{(Werner, 2017)},
w~pliku \texttt{wisiepk.cls} należy zmienić \texttt{style=numeric-comp} na
\texttt{style=authoryear} i~używać polecenia \verb|\parencite|.

\section{Przykłady opisów bibliograficznych}
\label{sec:przyklady-bib}

Tabela~\ref{tab:typy-bib} zestawia najczęściej cytowane rodzaje źródeł
wraz z~typem wpisu w~pliku \texttt{literatura.bib} i~przykładem
z~wykazu literatury tej formatki.

\begin{table}[htbp]
  \centering
  \caption{Rodzaje źródeł i~odpowiadające im typy wpisów}
  \label{tab:typy-bib}
  \begin{tabularx}{\linewidth}{@{}>{\raggedright\arraybackslash}X
                                   >{\ttfamily}l c@{}}
    \toprule
    \textbf{Rodzaj źródła} & \textbf{\normalfont\bfseries Typ wpisu} & \textbf{Przykład} \\
    \midrule
    książka (podręcznik, monografia)       & @book          & \cite{pelech2009,duffie2013,metcalf2014} \\
    rozdział w~pracy zbiorowej             & @incollection  & \cite{nuorkivi2016} \\
    artykuł w~czasopiśmie                  & @article       & \cite{lund2014,werner2017} \\
    referat konferencyjny                  & @inproceedings & \cite{crawley2005} \\
    raport, instrukcja, dokumentacja       & @report        & \cite{iea2022,epanet22} \\
    norma krajowa lub międzynarodowa       & @standard      & \cite{pn12831,pn6946,pn14001,iso52000} \\
    ustawa, rozporządzenie, dyrektywa      & @legislation   & \cite{ustawaPOS,rozpWT,epbd2024} \\
    strona internetowa, serwis             & @online        & \cite{pvgis,tlm,isap} \\
    zbiór danych                           & @dataset       & \cite{imgw} \\
    oprogramowanie                         & @software      & \cite{qgis} \\
    karta katalogowa producenta            & @manual        & \cite{katalog} \\
    praca dyplomowa                        & @thesis        & \cite{praca2024} \\
    \bottomrule
  \end{tabularx}
\end{table}

\subsection{Normy}

W~opisie normy podaje się jej pełne oznaczenie \emph{wraz z~datą
wydania} (np.\ PN-EN 12831-1:2017-08), tytuł, wydawcę i~rok. Cytuje się
wersję faktycznie wykorzystaną w~obliczeniach, a~jej aktualność warto
sprawdzić w~katalogu Polskiego Komitetu Normalizacyjnego. Przykładowy
wpis pokazano na listingu~\ref{lst:bib-norma}.

\begin{lstlisting}[language={}, caption={Wpis normy w~pliku
  literatura.bib}, label={lst:bib-norma}]
@standard{pn12831,
  number       = {PN-EN 12831-1:2017-08},
  title        = {Charakterystyka energetyczna budynków -- Metoda
                  obliczania projektowego obciążenia cieplnego --
                  Część 1: Projektowe obciążenie cieplne, Moduł M3-3},
  organization = {Polski Komitet Normalizacyjny},
  location     = {Warszawa},
  year         = {2017},
}
\end{lstlisting}

\subsection{Akty prawne}

Ustawy i~rozporządzenia cytuje się według \emph{aktualnego tekstu
jednolitego} ogłoszonego w~Dzienniku Ustaw, z~dopiskiem „z~późn.~zm.”, jeśli
po jego ogłoszeniu wprowadzono zmiany. Stan prawny i~numer tekstu
jednolitego najłatwiej sprawdzić w~Internetowym Systemie Aktów Prawnych
\cite{isap}. Akty prawa Unii Europejskiej cytuje się według Dziennika
Urzędowego UE \cite{epbd2024}.

\subsection{Źródła internetowe}

Przy stronach internetowych obowiązkowo podaje się adres URL oraz datę
dostępu (pole \texttt{urldate}), ponieważ treść strony może się zmienić
(listing~\ref{lst:bib-www}). Należy korzystać przede wszystkim ze
źródeł instytucjonalnych (urzędy, instytuty badawcze, organizacje
międzynarodowe). Encyklopedie redagowane przez użytkowników, fora i~blogi
nie są wiarygodnym źródłem w~pracy dyplomowej.

\begin{lstlisting}[language={}, caption={Wpis strony internetowej
  w~pliku literatura.bib}, label={lst:bib-www}]
@online{pvgis,
  author       = {{European Commission, Joint Research Centre}},
  title        = {Photovoltaic Geographical Information System (PVGIS)},
  url          = {https://joint-research-centre.ec.europa.eu/...},
  urldate      = {2026-10-09},
}
\end{lstlisting}

Nazwę instytucji występującej w~roli autora ujmuje się w~podwójne
nawiasy klamrowe, aby nie została potraktowana jak imię i~nazwisko.
Datę dostępu zapisuje się w~formacie RRRR-MM-DD -- w~wykazie literatury
zostanie wydrukowana w~formie polskiej.
